\section{An alternating tensor-orbit obstruction with arbitrary partners}
\label{sec:alternating-tensor-orbit}

\paragraph{Status and provenance.}
D-TD14--15, reserved question Q2ee, is a separate
\texttt{/root/imprimitive} derivation. It uses the accepted actual-cover
genus theorem HD-GC06 and the elementary trace and cyclic-orbit ideas
of HD-TC01, with the arbitrary-partner and full-lift interfaces proved
below. No external partner is discarded. The complete author whole-proof,
source-interface and code/output readings pass, and the exact supplement
passed its first execution. Root and independent whole-proof readings also pass without mathematical repair.
This does not change D-TD11--13 or HD-SS01--05.

Let a finite group \(J\) act on a complex vector space \(V\) of dimension
\(e\). Suppose there is an actual tensor decomposition
\[
 V=(W_1\otimes\cdots\otimes W_n)\otimes R,\qquad
 \dim W_i=a\ge2,\qquad \dim R\ge1,
 \tag{TO.1}
\]
such that \(J\) permutes the \(n\) embedded full matrix algebras
\(\operatorname{End}(W_i)\). The action on this set of algebras defines
an actual homomorphism
\[
                 x:J\longrightarrow S_n.
 \tag{TO.2}
\]
The external matrix algebra \(\operatorname{End}(R)\), as the
commutant of the product of the displayed factor algebras, is preserved.
Each operator is therefore a weighted permutation of the \(W_i\),
tensored with an invertible operator on \(R\). To see this directly,
first undo the permutation of the factor algebras; their remaining
automorphisms are inner. Undo those inner automorphisms as well.
The resulting operator centralizes their full tensor matrix algebra
and so acts only on \(R\). All choices of individual factor maps
are unique only up to scalars, which are allowed throughout.
This is the matrix-algebra interface of D-TD11--13, not an assertion
that the individual \(W_i\) are stable representations of \(J\).

\subsection*{D-TD14. Trace and eigenspace bounds for the actual lift}

\begin{lemma}[D-TD14]
Let \(T\) be any finite-order operator in the normalizer in (TO.1),
and let \(\pi\) be its induced factor permutation. For every integer \(j\),
\[
       \frac{|\operatorname{tr}(T^j)|}{e}
                   \le a^{-\operatorname{ind}(\pi^j)},
 \qquad
 \operatorname{ind}(\pi)=n-\#\operatorname{cycles}(\pi).
 \tag{TO.3}
\]
For every nonidentity \(\pi\), every scalar phase and every external
operator on \(R\), its maximum-eigenspace codimension satisfies
\[
                        \nu(T)\ge e/4.
 \tag{TO.4}
\]
If \(n\ge128\), then, on
\(\Omega=\{1,\ldots,n\}^3\) with repetitions allowed,
\[
      \frac{\nu(T)}e
       \ge\frac{\operatorname{ind}(\pi,\Omega)}{n^3}
                       -\frac1{n^6}.
 \tag{TO.5}
\]
No assertion that \(T^{\operatorname{ord}(\pi)}\) is scalar is used.
\end{lemma}

\begin{proof}
Write \(T=A\otimes B\), with \(A\) the weighted tensor permutation
on the first \(n\) factors. If \(T^L=1\), then
\(A^L=cI\), \(B^L=c^{-1}I\): a tensor product of two invertible
operators is scalar only if both are scalar. Choose
\(\alpha^L=c^{-1}\) and replace \(A,B\) by
\(\alpha A,\alpha^{-1}B\). Their tensor product is unchanged and
now both have \(L\)-th power one. Thus both are semisimple with
all eigenvalues on the unit circle, and \(L\) is a multiple of
\(\operatorname{ord}(\pi)\).

We prove the weighted trace bound without assuming an invariant
tensor-product metric. Around each factor cycle \(C\) of length
\(\ell_C\), conjugate the tensor identifications to put its cycle
product \(S_C\) on one factor and identities on the others.
The tensor trace identity gives
\[
                  \operatorname{tr}A=\prod_C\operatorname{tr}S_C.
 \tag{TO.6}
\]
Since \(A^L=1\), each \(S_C^{L/\ell_C}\) is a scalar \(c_CI\);
it occurs on \(\ell_C\) factors in \(A^L\), and
\(\prod_Cc_C^{\ell_C}=1\). Put
\(r_C=|c_C|^{\ell_C/L}\). Every eigenvalue of \(S_C\) has
absolute value \(r_C\), and \(\prod_Cr_C=1\). Consequently
\[
 |\operatorname{tr}A|
      \le\prod_C(ar_C)=a^{\#\operatorname{cycles}(\pi)}.
 \tag{TO.7}
\]
Apply the same argument to \(A^j\), whose factor permutation is
\(\pi^j\). Also \(|\operatorname{tr}B^j|\le\dim R\).
This proves (TO.3), including all powers with trivial factor
permutation and arbitrary behavior in the kernel of (TO.2).

For (TO.4), choose a factor cycle of length \(\ell\ge2\).
Its product \(S_C\) is diagonalizable by the preceding scalar-power
identity. In an eigenbasis the weighted cyclic operator is monomial
on the tensor-word basis. There are \(a\) constant words; every other
word orbit has length at least two. On a word orbit of length \(d\),
the weighted cyclic matrix has polynomial \(z^d-\beta\), with
\(\beta\ne0\), so a specified eigenvalue occurs at most once.
Its maximum eigenspace therefore has dimension at most
\((a^\ell+a)/2\). The support fraction on this cycle is at least
\((1-a^{1-\ell})/2\ge1/4\). Spectral convolution with every other
cycle and with \(B\) cannot increase the largest normalized
eigenvalue multiplicity. This proves the asserted bound for
every target eigenvalue, including all small cycle orders.

Here are elementary uniform constants for the cubic comparison.
For every integer \(n\ge128\),
\[
                          2^n>n^{12}.
 \tag{TO.8}
\]
It holds at \(128\), since \(128^{12}=2^{84}<2^{128}\).
For \(n\ge128\),
\[
 (1+1/n)^{12}
 \le\sum_{k=0}^{12}(12/n)^k
 <(1-12/n)^{-1}\le32/29<2,
 \tag{TO.9}
\]
so induction proves (TO.8).
If \(f\) letters of a permutation are fixed, its nontrivial
cycles give \(\operatorname{ind}(\pi)\ge(n-f)/2\).
For \(1\le f\le n\), consider
\[
 H(f)=\frac{n-f}{2}\log2+3\log(f/n).
 \tag{TO.10}
\]
It is concave on \([1,n]\), has \(H(n)=0\), and has
\(H(1)>0\) because (TO.8) implies \(2^{n-1}>n^6\).
It follows that \(H(f)\ge0\) throughout the interval.
Together with (TO.3) this proves
\[
 \frac{|\operatorname{tr}T|}{e}\le
 \begin{cases}
       (f/n)^3,&f>0,\\
       n^{-6},&f=0.
 \end{cases}
 \tag{TO.11}
\]
The second case uses \(2^{-n/2}<n^{-6}\), again from (TO.8).
The weak inequalities displayed are convenient and sufficient.
The same inequalities apply to every power of \(T\).

Let \(z\) be any eigenvalue of \(T\). Fourier inversion over its
full order \(L\), not just the order of \(\pi\), gives
\[
 \begin{split}
 \frac{\dim\ker(T-zI)}e
  &\le\frac1L\sum_{j=0}^{L-1}
                       \frac{|\operatorname{tr}T^j|}{e}\\
  &\le\frac1{Ln^3}\sum_{j=0}^{L-1} f(\pi^j)^3+n^{-6}\\
  &=\frac{\#(\Omega/\langle\pi\rangle)}{n^3}+n^{-6}.
 \end{split}
 \tag{TO.12}
\]
The last equality is Burnside's orbit formula; the sequence
\(\pi^j\) repeats equally often because its order divides \(L\).
This proves (TO.5). Powers in the kernel of the factor action
are bounded by one, rather than silently treated as scalar.
\end{proof}

\subsection*{D-TD15. The actual alternating-orbit exclusion}

\begin{theorem}[D-TD15]
Suppose \(C\longrightarrow\mathbf P^1\) is an actual connected
\(J\)-Galois cover whose generating product-one inertia
\(g_1,\ldots,g_r\) satisfies the sixfold coefficient budget
\[
             \sum_{i=1}^r\operatorname{rank}(1-V(g_i))=2e+6.
 \tag{TO.13}
\]
Assume the actual tensor structure (TO.1)--(TO.2) and
\[
                           x(J)=A_n\ \hbox{or}\ S_n.
 \tag{TO.14}
\]
Then \(n<10^6\). This conclusion allows any external partner,
any coefficient field, any scalar phases, and any number of
genuine projected inertia entries. It does not assume \(r\ge12\).
\end{theorem}

\begin{proof}
Suppose \(n\ge10^6\). Then \(e=a^n\dim R\ge2^n>24\).
Let \(r_0\) count the original inertia entries with \(x(g_i)\ne1\).
Every such entry costs at least \(e/4\) by (TO.4), so
\[
            r_0/4\le2+6/e<9/4,\qquad r_0\le8.
 \tag{TO.15}
\]
This count is a consequence, not a hypothesis. The remaining
entries have nonnegative moving rank and may be discarded.

The quotient \(C/\ker x\longrightarrow\mathbf P^1\) is an actual
connected \(A_n\)- or \(S_n\)-Galois cover. It has the projected
product-one generating tuple, with identity images omitted.
Its action on \(\Omega\) has five orbits: the diagonal, the three
two-equality patterns, and the ordered triples of distinct letters.
Write their connected quotient genera as
\(\gamma_1,\gamma_2,\gamma_2,\gamma_2,\gamma_3\).
Their exact Riemann--Hurwitz sum is
\[
 \sum_i\operatorname{ind}(x(g_i),\Omega)
        =2(n^3-5)+2(\gamma_1+3\gamma_2+\gamma_3).
 \tag{TO.16}
\]
The identity images make no contribution.

The accepted effective genus theorem HD-GC06 applies precisely to
this actual sphere cover, with \(n\ge10^6\), and gives
\(c>n^2/288\) for its unordered-triple quotient.
The ordered-distinct-triple curve maps to that curve with degree
six. For \(A_n\) as well as \(S_n\), all permutations of the three
selected letters are induced by the set stabilizer: an odd
permutation is compensated on two complement letters. Hence
Riemann--Hurwitz yields
\[
                  \gamma_3\ge6c-5>n^2/48-5.
 \tag{TO.17}
\]
The other genera are nonnegative. Thus the right side of
(TO.16) is strictly greater than \(2n^3+n^2/24-20\).

Apply (TO.5) to the \(r_0\) retained actual lifts and sum. Using
actual moving rank at least maximum-eigenspace codimension gives
\[
 \frac1e\sum_i\operatorname{rank}(1-V(g_i))
  >2+\frac1{24n}-\frac{20}{n^3}-\frac8{n^6}
  >2+\frac1{25n}.
 \tag{TO.18}
\]
For the last strict inequality it is enough that
\(n^5-12000n^3-4800>0\). This holds already for \(n\ge128\):
\(n^3(n^2-12000)\ge128^3(128^2-12000)>4800\).
Also \(2^n>150n\) for \(n\ge128\), by its base value and the
induction \(2(n)>n+1\). Since \(e\ge2^n\), the strict excess
in (TO.18) contributes more than six, contradicting (TO.13).
\end{proof}

\paragraph{Matched standard layer and exact boundary.}
In the single-standard setting \(V=\operatorname{Std}_n\otimes U\),
the theorem applies if \(U\) has an actual orbit of \(n\) equal
tensor factor algebras and its permutation homomorphism is the
same natural quotient \(x\). The standard module, and any other
factors, are then part of the unrestricted external \(R\).
This excludes that realization for \(n\ge10^6\) without proving
\(\mu_U=0\), changing the character field, or selecting a stable
individual factor from a permuted orbit.

More generally D-TD15 applies to any orbit of tensor factor
algebras whose entire actual permutation image is alternating
or symmetric in that degree. It does not apply merely because a
larger permutation image has an abstract alternating quotient.
No comparison between indices in two different permutation
actions is asserted. No cutoff on a different external standard
layer, no bound on all other tensor orbits, and no existence of
a source are inferred.

HD-TC02 had no unrestricted external partner and bounded tensor
count in terms of local degree. The current argument instead uses
the actual alternating quotient geometry and bounds this orbit
length independently of the local dimension and all partner
actions. No newly cited classification, character-table, or
asymptotic genus theorem enters the proof.
